\section{Original data and the equality theorems}
\label{sec:setting}

Throughout, \(n\ge3\), \(k=n-1\), and
\(\omega=|S^{n-1}|\) is the area of the unit round sphere. The Laplacian
is \(\Delta=\divg\nabla\), and symmetrization includes the factor \(1/2\).
Mean curvature is the tangential divergence of the specified normal.

\begin{definition}[One-ended initial-data exterior]
\label{def:data}
A one-ended initial-data exterior is a smooth connected orientable
\(n\)-manifold \(\Omega\) with nonempty compact smooth boundary \(S\),
a smooth Riemannian metric \(g\), and a smooth symmetric covariant
two-tensor \(K\), both smooth up to \(S\), satisfying the following
conditions.

The metric space \((\Omega,g)\), with \(S\) included, is complete.
There is a compact subset \(C\subset\Omega\) such that
\(\Omega\setminus C\) is a single coordinate end diffeomorphic to
\(\{x\in\R^n:|x|>R_0\}\). In these coordinates, for some
\[
 \frac{n-2}{2}<q<n-2,
\]
we have
\begin{equation}
 g_{ij}-\delta_{ij}=O_6(r^{-q}),\qquad
 K_{ij}=O_5(r^{-1-q}),\qquad r=|x|.
 \label{eq:setting:decay}
\end{equation}
Here \(O_j(r^{-a})\) includes the estimate
\(|\partial^\alpha T|\le C_\alpha r^{-a-|\alpha|}\) for every
\(|\alpha|\le j\). These are asymptotic bounds on otherwise smooth data.

Put
\begin{equation}
 \tau=\tr_gK,\qquad
 2\mu=R_g+\tau^2-|K|_g^2,\qquad
 J_i=\nabla^j(K_{ij}-\tau g_{ij}).
 \label{eq:setting:constraints}
\end{equation}
We assume the dominant energy condition and integrability:
\begin{equation}
 \mu\ge |J|_g,\qquad
 \mu,\ |J|_g\in L^1(\Omega,dV_g).
 \label{eq:setting:dec}
\end{equation}
The ADM limits
\begin{align}
 E&=\frac{1}{2k\omega}
 \lim_{R\to\infty}\int_{S_R}
   (\partial_jg_{ij}-\partial_ig_{jj})\,n_\delta^i\,dA_\delta,
 \label{eq:setting:energy}\\
 P_i&=\frac{1}{k\omega}
 \lim_{R\to\infty}\int_{S_R}
  (K_{ij}-\tau g_{ij})\,n_\delta^j\,dA_\delta
 \label{eq:setting:momentum}
\end{align}
exist and are finite. We impose the future-timelike condition
\begin{equation}
 E>|P|_\delta,\qquad m=(E^2-|P|_\delta^2)^{1/2}.
 \label{eq:setting:timelike}
\end{equation}

On every component of \(S\), the unit normal \(\nu\) points into
\(\Omega\). We require
\begin{equation}
 \theta_+(S)=H_S+\tr_SK\le0,\qquad
 H_S=\divg_S\nu,\qquad
 \tr_SK=(g^{ij}-\nu^i\nu^j)K_{ij}.
 \label{eq:setting:trapping}
\end{equation}
There is no assumption on the other null expansion.
\end{definition}

The compact-complement condition in Definition~\ref{def:data} is part of
our use of \emph{one-ended}. It excludes additional non-asymptotically-flat
ends. Merely specifying one asymptotically flat end among other
unrestricted ends would define a larger class, for which the compact
filling and global arguments below would need additional hypotheses.
The cosmological constant is zero. No spin, maximality, time-symmetry,
vacuum, stationary-development, or interior-topology assumption is made.

\begin{definition}[Full enclosing cuts]
\label{def:fullcuts}
A full enclosing cut is \(\Gamma=\partial D\), the entire intrinsic
manifold boundary of a connected smooth codimension-zero
submanifold-with-boundary \(D\) of \(\Omega\), such that \(D\) is closed
in \(\Omega\), its manifold interior lies in \(\operatorname{int}\Omega\),
and it contains the whole sufficiently distant coordinate end.
We require \(\Gamma\) to be compact, smooth, embedded, and two-sided.

All components of \(\partial D\) are counted, including any portion
coincident with \(S\). In particular, \(D=\Omega\) is admissible and
has full boundary \(S\). Define
\begin{equation}
 A_{\min}(S)=\inf_{\Gamma}\Area_g(\Gamma),
 \label{eq:setting:minarea}
\end{equation}
and assume \(A_{\min}(S)>0\).
\end{definition}

A cut need not be connected, minimal, or trapped. The areas in
\eqref{eq:setting:minarea} are always computed in the original metric
\(g\). We neither assume that the infimum is attained nor assign a
classical expansion to a singular minimizing boundary.

\begin{theorem}[Numerical input]\label{thm:numerical-input}
The numerical spacetime Penrose theorem \cite[Theorem~1.3]{CompanionNumerical} states that every initial-data exterior in Definition~\ref{def:data}, with
the positive full-cut infimum in Definition~\ref{def:fullcuts}, satisfies
\begin{equation}
 \boxed{\quad
 m\ge\frac12
 \left(\frac{A_{\min}(S)}{\omega}\right)^{\frac{n-2}{n-1}} .
 \quad}
 \label{eq:setting:main}
\end{equation}
This input is used for nearby weakly trapped data as well as for the
original pair. It imposes none of the additional equality hypotheses
introduced below.
\end{theorem}

The normalization is geometric: the Schwarzschild--Tangherlini metric
of mass \(M>0\) is
\begin{equation}
 \begin{gathered}
 \mathbf G_M=-F(r)\,dt^2+F(r)^{-1}\,dr^2+r^2\sigma_{n-1},\\
 F(r)=1-\frac{2M}{r^{n-2}},\qquad
 r>r_M:=(2M)^{1/(n-2)}.
 \end{gathered}
 \label{eq:setting:tangherlini}
\end{equation}
Here \(\sigma_{n-1}\) is the unit round metric. We use its regular
extension across the future horizon and at the bifurcation sphere.
For a spacelike hypersurface our sign convention is
\begin{equation}
 K(U,V)=\mathbf G_M(\nabla^{\mathbf G_M}_U n_{\mathrm{future}},V).
 \label{eq:setting:Ksign}
\end{equation}

\begin{definition}[Rigidity class]
\label{def:setting:rigidity}
An exterior in Definition~\ref{def:data} belongs to the rigidity class
if \(S\) is connected, \(\theta_+(S)=0\), every full cut has area at
least \(A=\Area_g(S)>0\), and no compact smooth embedded two-sided
enclosing hypersurface entirely in \(\operatorname{int}\Omega\),
possibly disconnected, has \(\theta_+\le0\) everywhere with normal
toward the end.
\end{definition}

Thus \(A_{\min}(S)=A\) in this class. Outermostness is imposed only for
smooth entirely interior enclosing hypersurfaces; its use on limiting
minimizers will be justified by the singular-sheet argument.

\begin{theorem}[Original-data rigidity]
\label{thm:rigidity}
Suppose an exterior in Definition~\ref{def:setting:rigidity} attains
equality in \eqref{eq:setting:main}. There is a global smooth spacelike
embedding of the entire \(\Omega\), including \(S\), into the regular
extension of one exterior of \(\mathbf G_m\), inducing exactly the
original \(g\) and \(K\) with convention \eqref{eq:setting:Ksign}.

The image lies in the domain \(r>r_m\) together with its corresponding
future horizon and bifurcation sphere. The boundary maps to a full
smooth section of that future horizon or to its bifurcation sphere.
The embedding and future unit normal extend smoothly to \(S\) in
horizon-regular coordinates, and the chosen end approaches spatial
infinity. There is no conclusion about a region behind \(S\).
\end{theorem}

A full horizon section meets every null generator once; its induced
metric has the same area as the round horizon sphere. In statements
about an asymptotic slice, approaching spatial infinity means that
the ambient static radius tends to infinity along every sequence
escaping the given coordinate end.

\begin{theorem}[Converse and sharpness]
\label{thm:converse}
Let \(M>0\). Every smooth spacelike exterior hypersurface in the regular
extension of \(\mathbf G_M\), with its interior in \(r>r_M\), whose
induced data satisfy
Definitions~\ref{def:data}, \ref{def:fullcuts}, and
\ref{def:setting:rigidity}, whose boundary is a full smooth section
of the corresponding future horizon or its bifurcation sphere,
and whose end approaches spatial infinity, has
\begin{equation}
 (E^2-|P|_\delta^2)^{1/2}=M,\qquad
 A=\omega(2M)^{\frac{n-1}{n-2}}.
 \label{eq:setting:converse}
\end{equation}
It therefore attains equality in \eqref{eq:setting:main}.
In every dimension there are admissible static examples and
admissible examples with \(K\not\equiv0\).
\end{theorem}

Theorems~\ref{thm:rigidity} and \ref{thm:converse} permit
non-time-symmetric and asymptotically boosted slices. The invariant
mass assertion in the converse refers to the original slice's ADM
integrals \eqref{eq:setting:energy}--\eqref{eq:setting:momentum};
it is not inferred only from an auxiliary Riemannian metric.


\subsection{The strict direction used in the variation}

The numerical inequality applies on a constraint set with boundary:
the dominant energy condition and the expansion condition may both be
saturated. The following analytic input supplies a direction into its
interior. Its proof is independent of the numerical inequality.

\begin{lemma}[Strict conformal direction]\label{lem:strict-input}
For the smooth exterior and decay in Definition~\ref{def:data}, choose
$0<\beta<\min(1,q)$ and a smooth positive extension $\varrho$ of the
coordinate radius. Let $B\geq|K|$ be smooth and nonnegative, with
$B=O_5(r^{-1-q})$. There is a unique smooth positive solution tending
to zero at infinity of
\[
 -\Delta_g\phi=B\sqrt{|d\phi|_g^2+\varrho^{-2k}}
                    +\varrho^{-n-\beta},\qquad
 \partial_\nu\phi=-1\quad\hbox{on }S.
\]
It satisfies $\phi=O_6(r^{2-n})$, and
$(-\Delta_g\phi)/r^{-n-\beta}\to1$. Its flux is finite and equals
\[
 f_\phi:=\lim_{R\to\infty}\int_{S_R}\partial_{\nu_g}\phi\,dA_g
 =-\Area(S)-\int_\Omega
 \left[B\sqrt{|d\phi|^2+\varrho^{-2k}}+\varrho^{-n-\beta}\right]dV_g<0.
\]
For sufficiently small $\delta>0$, the pair
$(e^{2\delta\phi}g,e^{\delta\phi}K)$ satisfies a strict global
weighted dominant energy inequality and has strictly negative future
expansion. Its energy is $E-\delta f_\phi/\omega$, its momentum is
$P$, and its metric is at least $g$ and converges uniformly relative
to $g$ as $\delta\downarrow0$.
\end{lemma}

This is the strict-direction lemma in the numerical companion
\cite[Lemma~3.5]{CompanionNumerical}; its
complete exterior proof accompanies that numerical theorem. Below we
also compute the precise active-ray derivative, since its positivity
is the part needed for separation.

\subsection{How the dimension-dependent conclusions fit together}

Theorem~\ref{thm:rigidity} concerns a one-ended exterior with strong
differentiated decay and requires no ambient extension behind its
boundary. There are further conclusions for weaker decay in spatial
dimensions three and four. Their precise hypotheses and statements
appear in Theorems~\ref{q3:intro:rigidity},
\ref{q3:intro:connected-rigidity}, and~\ref{q4:thm:rigidity}.

The three-dimensional statements assume a complete boundaryless
ambient initial-data manifold with finitely many asymptotically flat
ends. The finite-component version forbids weakly trapped enclosing
cuts even when they partly coincide with the given horizon; it forces
one spherical component. With a connected boundary, the weaker
wholly-interior outermostness condition suffices. These weak-decay
statements are forward rigidity results. The four-dimensional
statement is an exterior theorem, with a connected marginal
outer-area-minimizing horizon and wholly-interior outermostness; it
also includes its horizon-regular converse.

The distinction matters for the proof. The all-dimensional argument
excludes possible interior null ends by a circle extension before
forming a compact-core double. The low-dimensional classification
arguments permit such ends while applying their positive-mass inputs.
We will therefore match regularity and completion separately in each
branch, after localizing the original-data variational equations.
